\begin{pastexam}[label={px:nchu-114-dsa-23}]{中興資工所甲組 114 學年度 計算題第 3 題}
（計算）

(5\%) You are given the sequence of integers 15, 10, 20, 8, 12, 17, 25, 19 to insert into an initially empty binary search tree (BST). After constructing the BST:
\begin{enumerate}[label=\Roman*.]
\item Delete the node with value 20 and replace it with its in-order successor.
\item Perform a post-order traversal of the updated BST.
\end{enumerate}
What is the resulting post-order traversal of the BST after the deletion?
\end{pastexam}

\begin{pastsoln}
答案：$8, 12, 10, 19, 17, 25, 15$。

依序插入，每個鍵值比較過的節點（資結書第~9.2~節的寫法）：
\begin{center}
\begin{tabular}{cll}
\toprule
插入 & 比較過的節點 & 接在 \\
\midrule
15 & （空樹） & 成為根 \\
10 & 15 & 15 的左邊 \\
20 & 15 & 15 的右邊 \\
8 & 15, 10 & 10 的左邊 \\
12 & 15, 10 & 10 的右邊 \\
17 & 15, 20 & 20 的左邊 \\
25 & 15, 20 & 20 的右邊 \\
19 & 15, 20, 17 & 17 的右邊 \\
\bottomrule
\end{tabular}
\end{center}

得到下圖 (a)。 刪除 $20$：它有兩個子節點，中序後繼是右子樹 $\{25\}$ 的最小值 $25$。 $25$ 就是 $20$ 的
右子節點，沒有左子節點，直接坐到 $20$ 的位置，接手左子樹 $17(-,19)$，右邊是空的，得到圖 (b)。

\begin{center}
\begin{minipage}{0.48\linewidth}\centering
\begin{tikzpicture}[level distance=9mm, edge from parent/.style={draw=black!55, thick},
  every node/.style={circle, draw=black!65, minimum size=6mm, inner sep=0pt, font=\footnotesize},
  level 1/.style={sibling distance=26mm}, level 2/.style={sibling distance=13mm}]
\node {15}
  child {node {10} child {node {8}} child {node {12}}}
  child {node[draw=memStroke, thick, fill=memFill] {20}
    child {node {17} child[missing] child {node {19}}}
    child {node[draw=optStroke, thick, fill=optFill] {25}}};
\end{tikzpicture}

{\footnotesize (a) 插入完成；橘色是要刪的 $20$，粉紅色是後繼}
\end{minipage}\hfill
\begin{minipage}{0.48\linewidth}\centering
\begin{tikzpicture}[level distance=9mm, edge from parent/.style={draw=black!55, thick},
  every node/.style={circle, draw=black!65, minimum size=6mm, inner sep=0pt, font=\footnotesize},
  level 1/.style={sibling distance=26mm}, level 2/.style={sibling distance=13mm}]
\node {15}
  child {node {10} child {node {8}} child {node {12}}}
  child {node[draw=optStroke, thick, fill=optFill] {25}
    child {node {17} child[missing] child {node {19}}}
    child[missing]};
\end{tikzpicture}

{\footnotesize (b) 刪除 $20$ 之後}
\end{minipage}
\end{center}

後序走訪（先左、再右、最後自己，資結書第~7.3~節）：左子樹 $10(8,12)$ 給出 $8, 12, 10$；右子樹
$25(17(-,19),-)$ 給出 $19, 17, 25$；最後是根 $15$。 答案：
\[
8,\ 12,\ 10,\ 19,\ 17,\ 25,\ 15.
\]
檢查：中序走訪 $8, 10, 12, 15, 17, 19, 25$ 仍然遞增，而且少了 $20$。
\end{pastsoln}
